\documentclass[leqno, a4paper, 12pt]{jsarticle}
\usepackage{amssymb}
\usepackage{amsthm}
\usepackage{amsmath}
\usepackage{comment}
\usepackage{url}

\usepackage{enumitem}
%\usepackage{showkeys}


%定理環境 thmはあまり使わずpropをなるべく使う。
%主張は基本的に証明の中で使い、そのため番号を付けない。
%注意環境を付けてみた。
\newtheorem{thm}{定理}
\newtheorem{prop}[thm]{命題}
\newtheorem{cor}[thm]{系}
\newtheorem{lem}[thm]{補題}
\newtheorem{df}[thm]{定義}
\newtheorem{exam}[thm]{例}
\newtheorem{fact}[thm]{事実}
\newtheorem*{claim}{主張}
\newtheorem*{cau}{注意}
\newtheorem*{rmk}{注意}
\renewcommand{\proofname}{証明}
%矢印たち。
%\toは\rightarrowと同じ。\getsは\leftarrowと同じ。同値は\iff
\newcommand{\To}{\Longrightarrow}
\newcommand{\Gets}{\LongLeftarrow}
%黒板太字
\newcommand{\nn}{\mathbb{N}}
\newcommand{\zz}{\mathbb{Z}}
\newcommand{\qq}{\mathbb{Q}}
\newcommand{\rr}{\mathbb{R}}
\newcommand{\cc}{\mathbb{C}}
%無限大
\newcommand{\mgn}{\infty}
%差集合は\setminusで統一する。等号の否定は\neq
%真部分集合は\subsetneqq　偏微分は\partial
%ディスプレイスタイル。\dfracでディスプレイ分数
\newcommand{\dpst}{\displaystyle}
\newcommand{\card}{\mathrm{card}}


\newcommand{\cl}{\mathrm{CL}}
\newcommand{\nai}{\mathrm{Int}}
\title{論文メモ
\large{\\--- ブラウン・モース・サードの定理 ---}
}
\author{ナンブ　キトラ}
\date{\today}
\begin{document}
\maketitle

本棚を整理したいので，
溜まっている論文のメモを書いて未来への手紙とすることにする．


ここではときとしてブラウン・モース・サードの定理としても呼ばれるサードの定理についての論文を色々紹介する．


\section{サードの定理まで}
おそらくサードの定理の雛形になるのは
関数の従属関係に関する研究であろう．
とりあえずこの方面の研究は
1926年のクノップとシュミットの共著である
\cite{knopp1926funktionaldeterminanten}
にまで遡れると思う．
この論文は日本語による解説もある
\cite{1927227}．
しかしながらそれ以前から
ヤコビ行列式と関数それ自体の従属関係
については研究があったようである．

1929年の論文
\cite{ost1929}
では任意の閉集合が
何かしらの関数の臨界点(critical point)
全体と一致するようにできることが書かれているらしい．
この解説は
\cite{MR1545896}の注釈３による．

1929年の論文
\cite{MR1561794}
も臨界点の話をしているがよくわからず．

1929年の論文
\cite{MR1561686}
は
モースの論文である．

1935年のBrownの
論文
\cite{MR1501816}
の
Lemma 3. I. 
はサードの定理の原型だと思われる．

1939年のモースの論文
\cite{MR1503449}
では，
値域が$\rr$
の場合のサードの定理の主張が書かれており，
現在よりも大きな微分可能性を仮定した上で
それが成り立つことが述べられている．
つまり$G$を$\rr^{n}$の開集合として
$f\colon G\to \rr$
を
$r$
階微分可能関数とするとき$r$が十分大きければ
$f$の臨界値(critical value)
の値が測度ゼロであることが知られていいたことを述べている．
どれくらい大きければいいのかというとサードの未発表論文
においては$r\le n+(n-3)^{2}/16$
であれば良いということが\cite{MR1503449}の冒頭で述べられている．
この論文の主定理としては上で述べた
値域が
$\rr$
の場合には$r=m$のときに
も上の主張が正しいことを証明している．

1942年の
サードの論文
\cite{MR0007523}
では
今現在の知られている形でのサードの定理が証明されている．
つまり
$G$を$\rr^{m}$の開集合として
$f\colon G\to \rr^{n}$
を$C^{r}$級の関数としたとき，
もしも$r\ge \max\{m-n+1, 1\}$
ならば
その臨界値の測度がゼロであることを
証明している．
この論文では
モースの論文\cite{MR1503449}
での補題が証明に用いられていることに注意しよう．
ここら辺の事情がモース・サードの定理と呼ばれる所以だと思われる．

サードの
1958年の
論文
\cite{MR0100063}
では，サードの定理が本質的に
ハウスドルフ次元にまつわる現象であることが見抜かれている．例えば$f$の定義域内の集合$A$について
$f$のrankが$<s$となるならば$f(A)$
の$s$次元ハウスドルフ測度は
$0$になることなどが証明されている．

サードの
1965年の論文
\cite{MR0173748}
では臨界点の像の話が
バナッハ多様体の話へと拡張されている．
この論文には修正
\cite{MR0180649}
があることに注意しよう．

\section{サードの定理から}

1966年のFederer の論文
\cite{MR0196033}
は１ページしかない．
そこでサードの定理の変種の話をしているのだが，
証明は有名なFedererのGeometric Measure Theory
に入っていることを述べている．
論文というよりアナウンスというであろう．

1967年の論文
\cite{dubovitskii1967sets}
でも臨界値のハウスドルフ測度の話をしている．

1967年の論文
\cite{MR0220884}
では
普通にfunctional dependenceの話をしている．
多分この話題について知りたければこの論文を読むと良いのではないだろうか．


1968年の論文
\cite{pohovzaev1968set}
ではバナッハ空間上で定義されたフレドホルム作用素についての臨界値のハウスドルフ次元の話をしている．
こういうフレドホルム作用素の話でサードの定理をやるのは
Smaleがやり始めたらしい(An infinite version of Sard's theorem)．


1971年の論文
\cite{MR0276988}
ではサードの定理の
構成的証明(constructive proof)
を与えているらしい．

1983年の論文
\cite{MR0716263}
はなぜだか知らないが，
すごく面白そうな気がする．
この論文を読むと何かいいことがあるような気がするが
あんまり読めていない．ごめんね．

1986年の論文
\cite{MR0837817}
ではなんか色々やっている．
多分サードの定理の証明読みたい時はこれを読むといいんじゃないかな．

1987年の
論文
\cite{MR0925986}
では$C^{\infty}$関数に関するサードの定理に対して
簡略化した関数を与えている．
おそらく大抵の人が目にするサードの定理の証明は
この論文が元ネタなのではないだろうか？
違う可能性もある．


1991年の論文
\cite{MR1142720}
では関数の変数を二つの部分に分けて片側では
$C^{k}$でもう片側では
$C^{\omega}$
のときのサードの定理を証明している．


1992年の論文
\cite{MR1074748}
ではヘルダー連続性を混ぜた話をしている．
1993年の論文
\cite{MR1194950}
はその続きである．

1993年の論文
\cite{MR1112486}
では
$C^{p, \alpha}$
級関数についての
サードの定理が証明されている．
$C^{p, \alpha}$
というのは
$p$
階微分が局所
$\alpha$
次ヘルダーということである．


1994年の論文
\cite{MR1294335}
では
微分可能性とZygmund 関数を混ぜた関数についての
サードの定理を証明している．
$C^{k+\textrm{Zygmund}}$級関数の話ということ．
Zygmundはある種の微分可能性である．


1996年の論文
\cite{MR1390652}
では実数直線のコンパクト部分集合が
臨界値になるための条件を調べているようだ．
面白い論文だと思う．

2001年の論文
\cite{CGTT}
では
$C^{k, \alpha}$
級関数の臨界値のハウスドルフ次元の
評価を行なっているようだ．


2004年の論文
\cite{jiang2004general}
では
臨界値のより詳しいハウスドルフ次元の評価をしているようだ．


2006年の論文
\cite{Siberian}
では
$\Omega\subset \rr^{2}$
を開集合として
$v\colon\Omega\to \rr$
が$0\not \in \cl\nai Dv(\Omega)$
を満たすならば
$v$の臨界値の測度がゼロになることを証明している．
これが成り立つのはなんか
すごいような気がする．

2009年の論文
\cite{MR2515490}
では
curve selection lemma
というものを使って
サードの定理の証明を簡略化している．
ただしこの
curve selection lemma 
というのは
o-minimal structureの理解がないと証明できなさそうである．

2015年の論文
\cite{MR3859213}
では関数の微分に適切な条件を課すと
サードの定理の仮定を少し緩められるという
話をしているようだ．


2019年の論文
\cite{MR4060490}
は
論文
\cite{MR2515490}
の発展である．
Bochnak--{\L}ojasiewicz inequality 
の不等式というのがあってそれの一般化と見做せる
いい不等式を使ってサードの定理を証明している．
Bochnak--{\L}ojasiewicz inequality はcurve 
selection lemma 
というか
o-minimal 
というか
semi-analytic sets 
の話なので
なんかすごい難しい．


2021年の論文
\cite{MR4340761}
はヘルダーを混ぜたサードの定理の話だと思う．
多分これを読むといいんじゃないかな．


\section{サードの定理にまつわる反例}
微分可能性が低い場合にはサードの定理が成り立たなかったりするがそのような例を
構成した論文を紹介する．

サードの定理にまつわる最初の反例は
Whitney
の論文
\cite{MR1545896}
である．
Whitneyは自身の証明した有名な
Whitneyの拡張定理を用いて
$C^{1}$級関数$f\colon \rr^{2}\to \rr$
であってその臨界値が非退化な区間を含むものを
構成している．
注意すべきなのはサードの定理が証明される以前に
Whitneyは現在のサードの定理の
仮定を
ギリギリ満たさない反例を構成しているということである．


1942年の論文
\cite{MR0007929}
では
連続で単調増加な
特異関数
を構成しているが,
これは多分サードの定理に関係ないと思う．
紛れていた．


1965年の論文
\cite{MR0182024}
では無限次元多様体の間の写像に関する
サードの定理の反例の話のようだ．


1979年の論文
\cite{MR0600614}
では，$R_{n}$を
$\rr^{n}$の$n$次元立方体としたときに
$C^{1}$級関数
$f\colon R_{3}\to R_{2}$
であって，
全射でなおかつ
至る所のrankが$\le 1$
となるような例が構成されている．
これはハーシュの「微分トポロジー」
の演習問題に載っている未解決問題であった．
またこれはサードの定理の仮定も満たさないし，
結論も満たさない．



1985年の論文
\cite{MR0820057}
では，カントール集合$C$について
$C+C=\{\, x+y\mid x, y\in C\, \}=[0, 2]$
という現象を用いてサードの定理の仮定を満たさないし
結論も満たさない例を構成している．


1989年の論文
\cite{MR0969524}
では
quasi-arcが
関数の臨界点になるための十分条件が考察されている．



1991年の論文
\cite{MR1142720}
では関数の変数を二つの部分に分けて片側では
$C^{k}$でもう片側では
$C^{\omega}$
のときのサードの定理を証明しているということを上で紹介したが，
論文
\cite{MR1142725}
では片側を
$C^{\infty}$にすると
Pughの結果が成り立たないことを示しているらしい．

1993年の
論文
\cite{MR1217168}
では任意の可分バナッハ空間$E$に対して
全射$C^{1}$級関数$f\colon E\to \rr^{2}$
であってその微分のランクは至るところ
$\le1$になっているものを構成している．
これはKaufmanの論文
\cite{MR0600614}
の発展と見做せる．

1999年の論文
\cite{MR1694403}
では
Whitney arcというものを構成している．
これはおそらくWhitneyの論文
\cite{MR1545896}に出てくるようなarcのことだと思われる．

2003年の論文
\cite{MR1991757}
では
Assouadの埋め込み定理を用いて
Whiney型の例を構成している．



2008年の
論文
\cite{MR2399495}
もWhitneyの例の話をしている．


2009年の論文
\cite{MR2581221}
は
Whitneyの例を受け，
連結集合が臨界点になるような特徴付けを発見したらしい．
すごい．


2011年の論文
\cite{MR2894520}
では
滑らかなある曲線上で
接微分が消えるような
定数でない連続関数を構成しているようだ．


2020年の論文
\cite{MR4172678}
は面白いと思う．
$\mathbb{S}^{n+1}$から
$\mathbb{S}^{n}$への写像について
ホモトピー的な考察と
サード的な考察を行なっており大変興味深いと思う．
ちょっと細かいところはわかんなかった．

\section{サードの定理の系}

1963年の論文
\cite{MR0145502}
では
サードの定理を用いて
$n$次元単体$E$について
その境界$\partial E$
は$E$のレトラクトではないことを証明している．
このことから
ブラウワーの不動点定理などを証明できる．

1975年の論文
\cite{MR0364578}
では$K\subset \rr^{n}$として
$f\colon K\to \rr^{n}$をリプシッツ写像
として$K$のルベーグ測度が有限としたときに
$C=\{\, y\in \rr^{n}\mid \card(f^{-1}(y))=\infty\, \}$
がルベーグ測度で測度$0$
になることを述べている．
$f$が$C^{1}$級のときはこのことはサードの定理の系である．

\section{サードの定理の変種}
可微分ではなく他のいろいろな種類の関数についての
サードの定理の類似が色々証明されているようだ．

\subsection{Delta-convex}
2008年の論文
\cite{MR2472482}
ではバナッハ空間の間のdelta-convex関数という
クラスについて，
定義域が
$\rr$の区間の場合にはその臨界値の
$1/2$ハウスドルフ測度はゼロになることを証明している．


\subsection{Borel}
論文\cite{MR1305213}
ではBorel写像に関するサードの定理のようなものを
証明している．
サードの定理の系として
$f\colon \rr^{m}\to \rr^{n}$をいい感じに滑らかな関数としたときに
$f^{-1}(y)$が大体の$y$について
$m-n$次元
可微分多様体になっているというものがある．
この論文では
適切な条件のもと，ボレル写像$f\colon \rr^{m}\to \rr^{n}$
について$f^{-1}(y)$がほとんどの$y$について
ハウスドルフ次元が$m-n$以下ということを示している．


\subsection{ソボレフ}

ソボレフ空間に属する関数についての
サードの定理に関する論文は
以下の通りである．
\cite{MR1871360}(大元の論文?), 
\cite{MR2183045},
\cite{MR2415054}(simple proof?), 
\cite{MR2923414}, 
\cite{MR3318511}. 

\subsection{リーマン多様体の距離関数}
論文
\cite{MR2128549}
はリーマン多様体の距離関数
に関するサードの定理の話である．

\subsection{よくわかんないやつ}
多分サードの定理に関係するが，
よくわからなかったやつ．
\cite{MR0154296}, 
\cite{MR2355374}. 


以上の論文以外にも本当にたくさんのサードの定理に関係する論文があるとは思うがおそらくこれらの論文の参考文献やこれらの論文を引用している論文から全て見つけ出せるのではないだろうか．

%READ!
%myplain is the same to amsplain style. 
\nocite{*}
\bibliographystyle{myplaindoidoi}
\bibliography{sard1.bib}



\end{document}